\section{Stage 2: the spectral core}\label{sec:core}
\sectionbrief
{one field $\hat u(t,x)$ on a uniform grid (one posterior sample), with a noise level $\sigma$; $x$ periodic or windowed.}
{to extract from $\hat u$ the factors of the symbol, the component of each factor, its weight and the quantities $\gamma_N$, $c_F$, $\rho$, together with a posterior probability on the factor structure.}
{the only deterministic stage, called as a black box on each sample. It implements the analytic theorems (factor decomposition, interaction term) on numerical data. Two versions exist: v1 (exact data) and v2 (structured fit, used with noise).}

\begin{symbols}
\begin{tabular}{@{}L{2.5cm}L{6.0cm}L{4.4cm}@{}}
$\hat u(t,\kappa)$ & Fourier transform of $u$ in $x$ (Sec.~\ref{sec:prelim}); on a grid, the FFT & computed from the sample\\
$n_x,\ dt$ & number of grid points in $x$; time step & the grid\\
$y_k,\ z_m$ & time series $y_k=\hat u(t_k,\kappa)$ at one fixed $\kappa$; and $z_m=e^{\lambda_mdt}$ & definitions of Sec.~\ref{sec:modes}\\
$A_m,\ \lambda_m$ & amplitude and rate of the $m$-th exponential at that $\kappa$ & unknowns of the pencil\\
$Y,\ Y_1,\ Y_2,\ V$ & Hankel matrix of $y$; $Y$ without last/first column; its leading right singular vectors & linear algebra of the pencil\\
$\mathrm{snr}_{\min}$, $\sigma$, $\sigma_{\mathrm{bin}}$ & minimum signal-to-noise ratio ($=3$); noise std of the field; noise std of one FFT bin, $\sigma/\sqrt{n_x}$ & chosen default; estimated\\
point $p=(\kappa_p,\lambda_p)$, $n$ & one extracted mode, and the number of them & output of Sec.~\ref{sec:modes}\\
$\sigma_p\ (=\sigma_\lambda)$ & uncertainty (std) of $\lambda_p$, from the Cram\'er--Rao bound & Eq.~\eqref{eq:crb}\\
$J,\ g_j$ & Jacobian matrix of the exponential model and $g_j(t)=e^{\lambda_jt}$ & used in Eq.~\eqref{eq:crb}\\
$d_j(p)$ & distance of point $p$ to the curve $Q_j=0$, in units of $\sigma_p$ & Thm.~\ref{thm:structfit}\\
$\ell(p)$ & the label (index of the factor) assigned to $p$ & Thm.~\ref{thm:structfit}\\
$\chi^2,\ s,\ n_{\mathrm{par}}$ & robust sum of squared distances; noise scale ($\ge1$); number of fitted coefficients & Eq.~\eqref{eq:bic}\\
$\mathrm{BIC}$ & Bayesian information criterion: fit plus a penalty per parameter & statistics\\
$u_j,\ w_j$ & component of factor $j$ and its weight & Eqs.~in Sec.~\ref{sec:core}\\
$a(\xi)=\sum_nc_n\xi^n$ & for $m=1$: $Q=\tau-a(\xi)$, so $a$ is the right-hand side of $u_t=a(\partial_x)u$; $c_n$ are the \emph{term coefficients} & theory (terms of the operator)\\
$a_n,\ T,\ \hat u_0$ & the $n$-th term $c_n(i\kappa)^n$; time span of the data; initial profile $\hat u(t_0,\kappa)$ & definitions\\
$\Eint,\ \rho$ & interaction term and its relative size & Eq.~\eqref{eq:rho}\\
$B_j,\ S$ & orthonormal basis of the sampled kernel of $Q_j$; the matrix $[B_1\cdots B_N]$ & Eq.~\eqref{eq:gammaN}\\
$\gamma_N,\ c_F,\ \delta(P)$ & separability numbers and the size of the ambiguity & explained in the text\\
\end{tabular}
\end{symbols}

\subsection{Modes: FFT in $x$ and a matrix pencil in $t$}\label{sec:modes}
For each wavenumber $\kappa$ of the FFT, the time series $y_k=\hat u(t_k,\kappa)$ is a sum of exponentials,
\begin{equation}\label{eq:expsum}
\hat u(t,\kappa)\;\approx\;\sum_{m}A_m(\kappa)\,e^{\lambda_m(\kappa)t},
\end{equation}
because each mode $e^{\lambda t+i\kappa x}$ contributes only to the bin $\kappa$. The rates $\lambda_m(\kappa)$ are found by the \emph{matrix pencil} (Hua--Sarkar), a linear-algebra form of Prony's method. The result is a cloud of points $(\kappa,\lambda,A)$: the data of Section~\ref{sec:prelim}.

\begin{explained}{the matrix pencil, step by step}
\textbf{Step 1: samples as powers.} With $dt$ the time step and $z_m=e^{\lambda_mdt}$, the samples are $y_k=\sum_mA_mz_m^k$. Finding the $\lambda_m$ means finding the $z_m$.

\medskip\textbf{Step 2: the Hankel matrix has rank $r$.} Build $Y_{ik}=y_{i+k}$. Each row is a combination of the $r$ geometric sequences $(z_m^{i+k})_k$, so $\operatorname{rank}Y=r$, the number of exponentials. \emph{Example}: $y_k=3(1/2)^k+(1/3)^k$ gives a Hankel matrix of rank $2$ for every size $\ge2$. The order is read from the singular values: those above $\text{tol}\cdot s_{\max}$ are signal, the rest noise.

\medskip\textbf{Step 3: shift invariance.} Let $Y_1$ be $Y$ without its last column and $Y_2$ without its first. Then $Y_2=Y_1\,\mathcal Z$ on the signal subspace, where $\mathcal Z$ has eigenvalues $z_m$. In practice one projects onto the $r$ leading right singular vectors $V$ and takes the eigenvalues of $(V_1)^{+}V_2$, with $V_1,V_2$ the rows of $V$ without the last, resp.\ first, one.

\medskip\textbf{Step 4: rates and amplitudes.} $\lambda_m=\log z_m/dt$. The amplitudes follow from a linear least-squares fit of $y$ on $e^{\lambda_mt_k}$.

\medskip\textbf{Step 5: noise.} A bin is used only if its amplitude exceeds $\mathrm{snr}_{\min}=3$ times the noise level of the bin ($\sigma/\sqrt{n_x}$). In v2 the order is chosen \emph{per bin} by BIC among $r_0,\dots,r_0+2$, where $r_0$ is the Hankel rank above noise: an exponential merged into another one (unresolved) gets its own pole when the data support it. The pencil error under noise has no theoretical bound (Hankel/Prony conditioning is open, \thy{T10}); this is where experiments, not theorems, decide.
\end{explained}

\subsection{How uncertain is each point}
The structured fit of the next subsection compares points to curves \emph{in units of each point's own uncertainty}. For a mode the uncertainty $\sigma_\lambda$ is the Cram\'er--Rao bound of its rate: with $g_j(t)=e^{\lambda_jt}$ and the Jacobian columns $\big(g_j,\ A_jt\,g_j\big)$ for \emph{all} exponentials in the same FFT bin,
\begin{equation}\label{eq:crb}
\Cov(\lambda_j)\;=\;\sigma_{\mathrm{bin}}^2\,\big[(J^{H}J)^{-1}\big]_{\lambda_j\lambda_j}.
\end{equation}
The bin-wide Fisher matrix matters: a first version that used a $2\times2$ matrix per exponential underestimated $\sigma_\lambda$ when two exponentials interfered, and the fit then overfitted. A floor $s\ge1$ on a global noise scale (Jeffreys prior) prevents exact data from being overfitted.

\subsection{From ``fit one $P$, then factor'' to a structured fit}\label{sec:structfit}
\textbf{Core v1} fits the lowest-degree monic $P$ vanishing on the points (weighted least squares on the monomials of \eqref{eq:monic}) and then factors it numerically by tracking the roots $\lambda_r(\xi)$ along a complex path and testing which groups of branches have polynomial elementary symmetric functions. It is exact on exact data and is kept as the test oracle.

It fails under noise. Experiment E1: with relative noise of only $10^{-4}$, Klein--Gordon and the telegraph equation were \emph{falsely} factored, transport plus diffusion failed from $10^{-4}$, the wave equation at $10^{-1}$. The reason is structural: factoring a polynomial with floating-point coefficients is ill-posed, and no data-free threshold exists \thy{T15} (CONV-T5: the discrepancy rule picks the right order, but the threshold cannot be estimated from the data).

\textbf{Core v2} drops the factoring step. Every candidate \emph{structure} $\mathcal M=((m_1,dx_1),\dots,(m_k,dx_k))$ is fitted directly on the points, and the structures are compared statistically.

\begin{theorem}{(Structured factor fit: algorithm) \thy{T15, T10} \impl}\label{thm:structfit}
Let the points be $p=(\lambda_p,\kappa_p)$ with uncertainties $\sigma_p$. For a structure $\mathcal M$ with $k$ factors, alternate until the labels stop changing:
\begin{enumerate}[nosep]
\item \textbf{assign} each point to the factor curve nearest in uncertainty units,
\[
\ell(p)=\argmin_j\ d_j(p),\qquad d_j(p)=\frac{1}{\sigma_p}\,\frac{|Q_j(\lambda_p,i\kappa_p)|}{|\partial_\tau Q_j(\lambda_p,i\kappa_p)|};
\]
\item \textbf{refit} every $Q_j$ by weighted least squares on its own points (the weights are $1/\sigma_p$).
\end{enumerate}
This is ``$k$-means with algebraic curves in place of centres'' (a $k$-varieties regression); 10 initialisations are tried (2 deterministic by sorting branches, 8 random) and at most 30 iterations. The fit is scored by the robust residual
\begin{equation}\label{eq:bic}
\chi^2=\sum_p\min\big(z_p^2,\ (3s)^2\big),\quad z_p=\min_jd_j(p),\quad
\mathrm{BIC}(\mathcal M)=2n\ln\!\Big(\max\Big(\tfrac{\chi^2}{2n},1\Big)\Big)+n_{\mathrm{par}}\ln(2n),
\end{equation}
where $n$ is the number of points, $s^2=\max\big(1,\operatorname{median}(z^2)/0.455\big)$ is a robust noise scale ($0.455$ is the median of $\chi^2_1$), and $n_{\mathrm{par}}=\sum_jm_j(dx_j+1)$ complex coefficients (counted twice unless the field is real). The posterior over structures is
\begin{equation}\label{eq:postM}
p(\mathcal M\mid\text{data})\ \propto\ e^{-\mathrm{BIC}(\mathcal M)/2}.
\end{equation}
\end{theorem}
\begin{intuition}
Equation~\eqref{eq:bic} is the profile log-likelihood of Gaussian residuals whose scale $s$ is unknown but cannot be smaller than the Cram\'er--Rao floor ($s\ge1$), plus the BIC penalty (a Laplace approximation of the evidence). The cap $(3s)^2$ lets a few points be \emph{outliers} (an unresolved exponential merged into another) instead of corrupting the fit. The candidate set has $85$ structures: up to $3$ factors, $m_j\in\{1,2\}$, $\sum m_j\le3$, $dx_j\le4$, in canonical order so that no structure appears twice.
\end{intuition}

\begin{explained}{the fit on a toy cloud, and what the posterior says}
\textbf{Step 1: the points.} For $P=(\tau+\xi)(\tau-0.3\xi^2)$ and $\kappa=1,\dots,6$ there are two points per $\kappa$: $\lambda=-i\kappa$ (transport) and $\lambda=-0.3\kappa^2$ (diffusion).

\medskip\textbf{Step 2: the structure $((1,1),(1,2))$.} Two factors, each of $\tau$-degree $1$, in $\xi$ of degrees $1$ and $2$. Initialisation by sorting the two branches at each $\kappa$ by $\operatorname{Im}\lambda$ already separates the line (imaginary) from the parabola (real). Refitting each curve by least squares gives $\tau+\xi$ and $\tau-0.3\xi^2$ up to noise, and $d_j(p)$ is of order $1$ for every point.

\medskip\textbf{Step 3: the wrong structures.} $((1,2),)$ with a \emph{single} factor $\tau-a_1\xi-a_2\xi^2$ cannot pass through a line and a parabola at once: $\delta$ is large for half the points, $\chi^2$ is large, BIC is large. The structure $((1,1),(1,2),(1,2))$ fits equally well but carries more parameters: the BIC penalty makes it less probable. Two structures that differ only by a coefficient equal to $0$ (for example $(1,2)$ with $c_2\approx0$ versus $(1,1)$) are the same equation, and the posterior splits mass between them: in experiments the top weight ranged from $0.45$ to $0.98$.

\medskip\textbf{Step 4: reading the answer.} A low top weight is \emph{information}, not failure: it signals an ambiguity between structures that the data cannot separate. In the aggregated output (Section~\ref{sec:agg}) the posteriors of all samples are pooled.
\end{explained}

\subsection{Repeated factors: the confluent test}\label{sec:confluent}
If the symbol has a factor $Q^2$, the field contains terms $(A+Bt)e^{\lambda t}$ instead of simple exponentials, and \emph{no finite test can exclude this from finite data} (OPF5B-T6). The core handles it as follows: in each FFT bin with two close poles ($|\Delta\lambda|<4/T$) it compares, by BIC, two simple exponentials with refined poles (8 real parameters) against one confluent term (6 parameters). If the confluent model wins in any bin the output is labelled \textsc{undecided} with the reason written out. Before this test both cores returned a confident simple factor for $(\tau+\xi)^2$ (posterior $0.96$).

\subsection{Components, weights, shared modes}
Once the factors are known, each mode is assigned to a factor (the one with the smallest $d_j$), and the component is the sum of its modes,
\[
u_j(t,x)=\sum_{p:\ \ell(p)=j}A_p\,e^{\lambda_p(t-t_0)+i\kappa_px}.
\]
\paragraph{The ambiguity $\delta(P)$.} Two factors $Q_1,Q_2$ may have \emph{common solutions}: functions annihilated by both. If $f$ is one of them, then moving $f$ from $u_1$ to $u_2$ (replacing $u_1,u_2$ by $u_1+f,\ u_2-f$) changes neither $Q_1u_1=0$, nor $Q_2u_2=0$, nor the sum $u$. So the split of $u$ into components is not unique. The theory shows that the freedom is a \emph{finite-dimensional} space of exponential-polynomials, of dimension $\delta(P)=\sum_{i<j}\dim\C[\tau,\xi]/(Q_i,Q_j)$, which counts (with multiplicity) the points $(\tau,\xi)$ where $Q_i$ and $Q_j$ vanish together, \emph{i.e.} the modes $(\kappa,\lambda)$ that lie on both curves. Example: wave factors $Q_1=\tau-\xi$, $Q_2=\tau+\xi$ vanish together only at $(0,0)$, so $\delta(P)=1$: the only common solutions are constants. The same-looking number $n_{\mathrm{shared}}$ in the code counts such shared modes in the data.

A mode lying on \emph{two} varieties at once ($d<3$ for both: the intersection responsible for the ambiguity $\delta(P)$ of the theory, \thy{T2}) is split equally between them. This equal split is the \emph{minimum-norm representative} of the finite ambiguity and is the posterior mean of the Bayesian layer of Section~\ref{sec:bayes}. The weight is
\begin{equation}\label{eq:wj}
w_j=\frac{\|u_j\|_{L^2}}{\sum_k\|u_k\|_{L^2}}.
\end{equation}

\subsection{Terms inside a factor and the interaction ratio}
For an evolution factor $\tau-a(\xi)$ with $a(\xi)=\sum_nc_n\xi^n$, the terms are read from the \emph{branch} $\lambda(\kappa)=a(i\kappa)=\sum_nc_n(i\kappa)^n$ by regression on monomials: this is the log-symbol regression of the analytic note, without having to unwrap phases. The many different $\kappa$ matter: reading terms is a Vandermonde inversion, whose condition number blows up like $h^{-d}$ when the spectral points are clustered at spacing $h$ \thy{T9}. So the whole band is used, and the number of distinct $\kappa$ per factor is reported.

With $a_n(\kappa)=c_n(i\kappa)^n$ and $T$ the time span, the interaction ratio is the closed form of the analytic theory,
\begin{equation}\label{eq:rho}
\widehat{\Eint}=\Big[\prod_ne^{Ta_n}-1-\sum_n\big(e^{Ta_n}-1\big)\Big]\hat u_0,\qquad
\rho=\frac{\|\Eint\|}{\|u(T)-u_0\|},
\end{equation}
at cost $O(N)$ per wavenumber. It vanishes for a single term. Because $\Eint$ is not continuous in $H^m$, $\rho$ is used only as a \emph{distribution over posterior samples}; a small positive value for single-term factors in experiments ($\bar\rho\approx0.002$ at noise $10^{-3}$, $0.029$ at $10^{-2}$) comes only from spurious noise-level terms.

\subsection{Separability: $\gamma_N$ and $c_F$}
For each factor take the kernel functions $e^{\lambda_r(\kappa)t+i\kappa x}$ over the band of wavenumbers present (all roots $\lambda_r$ of $Q_j(\cdot,i\kappa)$), orthonormalise them to a basis $B_j$ of the sampled kernel, and set $S=[B_1\cdots B_N]$. Then
\begin{equation}\label{eq:gammaN}
\gamma_N=\min\{\mu\in\operatorname{spec}(S^{H}S):\ \mu>10^{-8}\},\qquad \gamma_2=1-c_F,
\end{equation}
and the number of eigenvalues $\le10^{-8}$ counts the dimension of the shared directions (the $\delta(P)$ ambiguity \thy{T2}). This is $\gamma_N$ of the analytic theory evaluated on the sampled cell, from the full block Gram matrix and \emph{not} from pairwise $c_F$, which do not determine $\gamma_3$ \thy{T5}. The factor $\gamma_N^{-1/2}$ amplifies every error on the components, and no regularisation can beat it \thy{T6}: it is an \emph{error floor}, used as an error bar and as a stop criterion. It depends strongly on the cell: for transport plus diffusion $\gamma\approx5\cdot10^{-4}$ for $T=0.1$, $0.17$ for $T=2$, $0.73$ for $T=8$; it falls about $300\times$ when the spatial window shrinks from $2\pi$ to $0.5$.

\subsection{The output contract}
\begin{caution}[title=The core never says ``identified'']
A fit from finite data cannot be certified (no finite certificate exists, \thy{T17}). The core returns one of four labels: \textsc{fit-compatible} (every factor has at least $m(dx+1)+1$ modes and at least $dx+1$ distinct $\kappa$), \textsc{id-ker} (the same, but only relative to the declared class, because there are too few modes: a single mode per factor fits infinitely many laws, \thy{T13}), \textsc{misfit} (more than half of the points are outliers for every structure), and \textsc{undecided} (no modes above noise, too few modes for any structure, or a confluent mode wins).
\end{caution}

\subsection{What has been measured}
On gridded data with $10$ cases (transport, heat, advection--diffusion, wave, beam, Klein--Gordon, telegraph, Airy, Schr\"odinger, transport plus diffusion), core v2 recognised the right factor set in $58$ of $60$ (case, noise) combinations up to relative noise $10\%$ in the regression test, and with five seeds per level in the benchmark: $1.00$ at noise $0$, $10^{-3}$, $10^{-2}$ and $0.90$ at $5\cdot10^{-2}$. Core v1 reached $0.66$, $0.60$, $0.56$ at the noisy levels. The best sparse-regression baseline (weak-form SINDy) reached $1.00$, $0.90$, $0.62$. The repeated-factor case $(\tau+\xi)^2$ is now reported as \textsc{undecided}. Open: a mixture outlier likelihood with marginalised outlier fraction (a fixed cap $3s$ is used), transport plus diffusion at noise $\ge3\cdot10^{-2}$, and a false \textsc{undecided} on the beam at $10^{-1}$. The same ten cases tuned the defaults and were then used to measure them; no held-out equations have been tried.

\begin{explained}{a reader asks: what is the matrix pencil for, and did we not already build the exponentials?}
\textbf{Step 1: what stage 1 produced.} Only numbers: a generic smooth field $\hat u_s$ sampled on a grid of $(t,x)$. It contains no exponential model; a GP draw is not written as a sum of exponentials.

\medskip\textbf{Step 2: building the exponential model, here.} The FFT in $x$ splits the field into bins, one per wavenumber $\kappa$; each bin is a function of $t$ alone. The matrix pencil fits that function as a sum of exponentials $\sum_mA_me^{\lambda_mt}$ and returns the rates $\lambda_m(\kappa)$. So space is handled by Fourier (modes $e^{i\kappa x}$), time by exponential fitting, and together they give modes $e^{\lambda t+i\kappa x}$: the exponential model in \emph{both} variables.

\medskip\textbf{Step 3: not yet the polynomial.} The pencil does \emph{not} estimate the coefficients of the polynomial. It only produces the points $(\kappa,\lambda)$. The coefficients of the factors are fitted afterwards, by regression of curves through those points (Section~\ref{sec:structfit}).
\end{explained}

\begin{takeaway}
Stage 2 turns each posterior sample into: a posterior over factor structures, the factors with their coefficients, components and weights, $\gamma_N$ and $c_F$, $\rho$, and a label that refuses to claim identification.
\end{takeaway}
